\documentclass[
    12pt,
    oneside,
    a4paper,
    brazil
]{abntex2}

\usepackage{hyperref}
\usepackage[utf8]{inputenc}
\usepackage{mathptmx}
\usepackage[T1]{fontenc}
\usepackage{indentfirst}
\usepackage{enumitem}
\usepackage{setspace}
\usepackage[section]{placeins}
\usepackage{amsmath}

% Informações
\titulo{Trabalho Final - Análise de Algoritmos}
\autor{Caio Collino Troiano \\ Gabriel Rosa Leite Barroso \\ Lucas da Mata Guimarães \\ Thiago Pereira Lins da Silva \\ Vinicius Vianna Egydio de Carvalho}
\local{São Paulo}
\data{2026}

\instituicao{
  Centro Universitário Senac
  \par
  Ciência da Computação
}

\hypersetup{
	colorlinks=true,
	linkcolor=blue
}

\tipotrabalho{Documento Técnico}
\preambulo{Documento de resultados de pesquisa desenvolvido para a disciplina de Análise de Algoritmos.}

\begin{document}

\imprimircapa

\tableofcontents*
\cleardoublepage

\OnehalfSpacing

\chapter{Rotting Oranges}

\textbf{Problema:} \href{https://leetcode.com/problems/rotting-oranges/}{Rotting Oranges}

\textbf{Paradigma:} Grafos (Busca em Largura / BFS Multi-source)

\textbf{Implementação:} (inserir o link pra qnd tiver o git)

\textbf{Complexidade:} (inserir complexidade)

\section{Estratégia de Solução}

O problema foi resolvido modelando a matriz $m \times n$  como um \textbf{grafo não-direcionado}, onde:
\begin{itemize}
	\item{Cada célula da matriz nas coordenadas $(i, j)$ funciona como um vértice;}
	\item{As conexões diretas (cima, baixo, esquerda, direita) funcionam como as arestas.}
\end{itemize}

Como a podridão se espalha de forma uniforme e simultânea a partir de várias fontes, a técnica ideal é a Busca em Largura Multi-Source (BFS). O uso de uma fila é essencial nessa abordagem, pois garante o processamento nível a nível, o que é fundamental para descobrirmos o tempo mínimo de infecção considerando que já existem várias laranjas podres no início. O algoritmo funciona em três etapas:
\begin{enumerate}
	\item{\textbf{Mapeamento Inicial:} O algoritmo percorre a matriz para contar o número total de laranjas frescas e insere a posição de todas as laranjas podres iniciais na fila da BFS;}
	\item{\textbf{Propagação por Minuto:} A fila é processada nível por nível (onde cada nível completo representa 1 minuto). Para cada laranja podre retirada da fila, o algoritmo infeta os seus vizinhos frescos válidos, mudando o estado deles para podre e colocando-os na fila para o minuto seguinte;}
	\item{\textbf{Verificação Final:} Quando a fila fica vazia, se o contador de laranjas frescas ainda for maior que zero, significa que algumas laranjas ficaram isoladas no grafo, retornando -1. Caso contrário, retorna o total de minutos decorridos.}
\end{enumerate}

\chapter{Nearest Exit from Maze}

\textbf{Problema:} \href{https://leetcode.com/problems/nearest-exit-from-entrance-in-maze/description/}{Nearest Exit from Maze}

\textbf{Paradigma:} Grafos (Busca em Largura / BFS Multi-source)

\textbf{Implementação:} (inserir o link pra qnd tiver o git)

\textbf{Complexidade:} (inserir complexidade)

\section{Estratégia de Solução}

O problema foi resolvido modelando a matriz $m \times n$ (o labirinto) como um grafo não-direcionado e não-ponderado, onde:
\begin{itemize}
	\item{Cada célula livre (caractere '.') da matriz nas coordenadas $(i, j)$ funciona como um vértice;}
	\item{As conexões diretas (cima, baixo, esquerda, direita) entre células livres adjacentes funcionam como as arestas de peso unitário.}
\end{itemize}

Como o objetivo principal é encontrar o caminho mais curto (menor número de passos) de um ponto de origem até à borda do labirinto, a técnica ideal é a Busca em Largura (BFS). O uso de uma fila (queue) é essencial nesta abordagem, pois garante uma expansão em "ondas" (nível a nível). Isto é fundamental para o problema, pois assegura que a primeira célula de saída encontrada seja, garantidamente, a mais próxima. O algoritmo funciona em três etapas principais:
\begin{enumerate}
	\item{\textbf{Mapeamento Inicial:} O algoritmo recolhe as coordenadas da entrada (entrance) e insere-as na fila da BFS com uma contagem de zero passos. Para otimizar o espaço espacial (evitando criar uma matriz extra de visitados), a própria célula de entrada é imediatamente marcada no labirinto original como visitada (mudando o caractere para '+'). Note-se que a célula de entrada, mesmo que esteja na borda, não conta como saída válida;}
	\item{\textbf{Busca e Expansão (Passo a Passo):} A fila é processada vértice a vértice. Para cada célula retirada, o algoritmo calcula as posições dos seus quatro vizinhos. Se um vizinho for válido (estiver dentro dos limites e for um caminho livre '.'), o algoritmo marca-o como visitado ('+'). Imediatamente a seguir, verifica se essa célula vizinha se encontra numa das bordas da matriz (linha ou coluna igual a $0$, ou igual ao limite máximo). Se for uma borda, a saída mais próxima foi encontrada e o algoritmo retorna de imediato o número de passos atuais mais um. Caso contrário, coloca o vizinho na fila para a expansão do próximo nível;}
	\item{\textbf{Verificação Final:} Quando a fila de processamento fica completamente vazia, significa que todos os caminhos possíveis foram explorados. Se a execução chegar a este ponto sem ter retornado nenhum valor na etapa anterior, conclui-se que não existe nenhuma rota possível para uma saída, retornando assim $-1$.}
\end{enumerate}

\chapter{Ice Statues Festival}

\textbf{Problema:} \href{https://judge.beecrowd.com/en/problems/view/1034}{Ice Statues Festival}

\textbf{Paradigma:} Programação Dinâmica (Bottom-Up / Tabulação)

\textbf{Implementação:} (inserir o link pra qnd tiver o git)

\textbf{Complexidade:} (inserir complexidade)

\section{Estratégia de Solução}

O problema foi resolvido modelando a busca pela quantidade mínima de blocos de gelo para atingir o comprimento desejado $M$ como uma tabela de estados de otimização (idêntico ao problema clássico Coin Change), onde:
\begin{itemize}
	\item{Cada índice $i$ do vetor (tabela) funciona como um subproblema, representando um comprimento específico de bloco que se deseja formar;}
	\item{O valor armazenado em cada índice $i$ representa a quantidade mínima de blocos de gelo necessários para formar aquele comprimento;}
	\item{Como o problema exige encontrar uma resposta ótima (menor número de blocos) e as soluções para valores maiores dependem diretamente das soluções para valores menores (caracterizando \textbf{subestrutura ótima} e \textbf{sobreposição de subproblemas}), a técnica ideal é a Programação Dinâmica;}
	\item{O uso de uma tabela (vetor) construída de baixo para cima (bottom-up) é essencial nessa abordagem, pois armazena os resultados parciais e evita o recálculo redundante, garantindo eficiência mesmo para valores de $M$ de até $1.000.000$.}
\end{itemize}

O algoritmo funciona em \textbf{três etapas}:
\begin{itemize}
	\item{\textbf{Inicialização da Tabela:} O algoritmo cria um vetor \texttt{dp} de tamanho $M + 1$. O caso base é definido como \texttt{dp[0] = 0}, estabelecendo que são necessários zero blocos para obter um comprimento de tamanho zero. Todos os outros espaços são inicializados com um valor muito grande (infinito simulado);}
	\item{\textbf{Construção Bottom-Up:} Através de laços aninhados, o algoritmo percorre a tabela de $1$ até $M$. Para cada posição $i$, ele avalia todos os blocos disponíveis cujos comprimentos sejam menores ou iguais a $i$. Utilizando a equação de recorrência \texttt{dp[i] = min(dp[i], dp[i - bloco] + 1)}, o algoritmo garante que o valor salvo seja sempre a melhor escolha;}
	\item{\textbf{Verificação Final:} Após preencher toda a tabela, a resposta ótima para o caso de teste estará no último índice (\texttt{dp[M]}), bastando imprimi-la. Como o bloco de tamanho $1$ sempre está presente, há a garantia de que uma solução viável sempre existirá.}
\end{itemize}

\chapter{Route Change}

\textbf{Problema:} \href{https://judge.beecrowd.com/en/problems/view/1123}{Route Change}

\textbf{Paradigma:} Grafos (Algoritmo de Dijkstra para Caminho Mínimo com Restrições)

\textbf{Implementação:} (inserir o link pra qnd tiver o git)

\textbf{Complexidade:} (inserir complexidade)

\section{Estratégia de Solução}

O problema foi resolvido modelando o sistema de estradas do país como um \textbf{grafo valorado} e \textbf{não-direcionado}, onde:
\begin{itemize}
	\item{Cada cidade representa um vértice (identificado de $0$ a $N-1$);}
	\item{Cada estrada representa uma aresta com peso equivalente ao custo do pedágio;}
	\item{Como o objetivo é encontrar o custo mínimo para atingir a cidade destino ($C-1$) a partir de uma cidade de quebra ($K$), considerando regras específicas para a rota de serviço fixa ($0, 1, \dots, C-1$), a técnica ideal é o \textbf{Algoritmo de Dijkstra}.}
\end{itemize}

A restrição fundamental do problema diz que, sempre que o veículo interceptar qualquer cidade que pertença à rota de serviço, ele \textbf{obrigatoriamente deve seguir a rota sequencial de forma linear} até o destino final, não podendo mais sair dela. O algoritmo funciona em \textbf{três etapas}:
\begin{enumerate}
	\item{\textbf{Modelagem e Adaptação do Grafo:} O grafo é montado por meio de uma matriz de adjacência (ou lista de adjacência). Para simplificar a restrição da rota de serviço, as arestas que partem de qualquer cidade $u$ pertencente à rota ($0 \le u < C$) são modificadas: ela só possuirá uma única saída válida direcionada para a cidade seguinte da rota ($u + 1$), com o peso correspondente ao pedágio dessa rota linear. Qualquer outra estrada saindo de $u$ é ignorada, pois o veículo não pode desviar se já estiver na rota;}
	\item{\textbf{Processamento do Caminho Mínimo (Dijkstra):} Executa-se o algoritmo de Dijkstra partindo do vértice inicial $K$ (a cidade onde o veículo foi consertado). Utiliza-se um vetor de distâncias \texttt{dist} inicializado com infinito e uma fila de prioridades (ou busca linear do menor vértice não visitado devido aos limites moderados de $N \le 250$) para computar o menor custo cumulativo acumulado até cada vértice do grafo;}
	\item{\textbf{Verificação Final:} Quando o algoritmo termina de processar os caminhos, o valor mínimo do custo acumulado para alcançar de forma válida o destino estará contido na posição correspondente à última cidade da rota de serviço (\texttt{dist[C - 1]}), sendo impresso como resultado para cada caso de teste.}
\end{enumerate}

\chapter{Beautiful Arrangement}

\textbf{Problema:} \href{https://leetcode.com/problems/beautiful-arrangement/description/}{Beautiful Arrangement}

\textbf{Paradigma:} Backtracking

\textbf{Implementação:} \href{https://github.com/BccMaterial/AA-Trabalho/blob/main/c/beautifulArrangement.c}{./c/beautifulArrangement.c}

\textbf{Complexidade:} $O(n!)$

\section{Estratégia de Solução}

\end{document}
